\documentclass[main.tex]{subfiles}

\begin{document}

\section{Introduction}
\label{sec:introduction}



Brain and central nervous system (CNS) tumors impose a substantial clinical burden because diagnosis, treatment planning, and follow-up depend on imaging evidence that must be interpreted under anatomical, pathological, and acquisition-related variability. Although brain/CNS tumors are less frequent than other cancer types, they are associated with severe neurological consequences and considerable mortality, with GLOBOCAN 2022 reporting more than 320,000 new cases and nearly 250,000 deaths worldwide \cite{globocan2022braincns,globocan2022chile}. Magnetic resonance imaging (MRI) is a core modality for brain tumor assessment because it provides non-invasive anatomical and tissue-dependent information used in diagnosis, follow-up, and treatment planning. However, MRI appearance can vary across tumor subtype, acquisition protocol, contrast sequence, scanner or site conditions, and patient-specific factors, making robust visual classification a non-trivial problem \cite{dorfner2025review,harmonizationMRI}.

Deep learning has therefore become a major research direction in brain tumor MRI analysis, including classification, segmentation, multimodal diagnosis, explainability, and representation learning \cite{dorfner2025review,ketabi2025multimodal,xaiMedical}. In many medical-image classification settings, the amount of labeled data is limited relative to the capacity of modern visual architectures. For this reason, practical workflows often adapt pre-trained visual backbones through fine-tuning or feature extraction rather than learning the full representation from scratch \cite{yosinski2014transferable,bar2024frofa}. Regardless of the specific backbone, these workflows share a common downstream structure: an image is encoded by a trained model, and the resulting representation is then classified, compared, adapted, or reused \cite{bar2024frofa,li2025embeddings}.


A common and natural choice is to use the final model-derived representation. Standard CNN classifiers typically map the last convolutional stage to a pooled feature vector and then to a learned classification head \cite{he2016resnet,huang2017densenet,tan2019efficientnet}. Consequently, once a trained backbone is available, the deepest embedding is often treated as the image descriptor for downstream classification, either through the native softmax head or through a classical classifier applied to the extracted feature vector. Classifiers such as k-nearest neighbors (kNN), support vector machines, random forests, and gradient-boosted trees provide standard downstream alternatives once such an embedding space is available \cite{cover1967nearest,cortes1995svm,breiman2001random,chen2016xgboost}. This final-embedding strategy is simple, reproducible, and computationally convenient. However, it also imposes a strong representational assumption: the last embedding is treated as the only relevant feature view for the downstream decision rule.

This assumption is not required by the architecture. CNNs produce a hierarchy of internal representations during the same forward pass. Earlier and intermediate layers tend to preserve more local or generic visual structure, such as edges, boundaries, textures, and mid-level patterns, whereas deeper layers become increasingly specialized for the training objective \cite{zeiler2014visualizing,yosinski2014transferable}. Prior work on feature visualization and transferability supports the idea that the semantic and geometric properties of CNN representations vary across depth \cite{zeiler2014visualizing,yosinski2014transferable}. Intermediate-representation and multilayer reuse methods further suggest that task-relevant information can be distributed across the network hierarchy rather than being restricted to the final layer \cite{fradi2021multilayer,gomes2023fmlf,evci2022head2toe}. Thus, the final descriptor is a strong and necessary baseline, but it is not the only representation exposed by a trained CNN.

The central methodological question is therefore not whether internal CNN layers can be used, but how they should be used under a valid evaluation protocol. Selecting a single layer can discard complementary information distributed across depth. Conversely, blindly concatenating or combining all available layers can create a heterogeneous and redundant high-dimensional descriptor. This issue is particularly important for distance-based classifiers, because high-dimensional spaces can reduce the discriminative value of pairwise distances through distance concentration, while redundant or poorly scaled feature blocks can distort neighborhood structure \cite{beyer1997nearest}. The opposite extreme is also relevant: overly aggressive compression may remove local, textural, shape, or contrast cues that earlier and intermediate stages preserve \cite{zeiler2014visualizing,yosinski2014transferable}. For brain tumor MRI, where subtle appearance differences and acquisition variability can influence class separability, the useful downstream representation is therefore not determined only by network depth; it also depends on how internal views are selected, scaled, compared, and combined.

Multilayer fusion methods differ in where the combination occurs. Architectural fusion combines features during model training or adaptation, often by adding trainable modules, attention mechanisms, or multiscale branches \cite{dai2020aff,he2025hemf,zhang2025context}. Decision-level fusion applies separate classifiers and combines their predictions after classification. Representation-level fusion constructs a shared descriptor before the final decision rule is applied. Distance-level fusion provides a related but distinct formulation: instead of concatenating coordinates into a single vector, it preserves each internal layer as a separate metric view and combines the distances induced by those views. This formulation is attractive in a post-hoc setting because it allows internal representations to contribute to the classifier without retraining the CNN backbone, introducing a new deep fusion module, or constructing a single high-dimensional concatenated descriptor.

This paper proposes \fullmethod{} (\method{}), a post-hoc method for constructing a distance-based classifier from internal CNN embeddings under validation control. Given a fine-tuned CNN backbone, the method extracts candidate internal representations from layers at different depths. Spatial activations are converted into vector embeddings through global average pooling (GAP) when necessary, and each resulting layer-specific vector is treated as a candidate internal embedding. Each layer view is then L2-normalized independently. The method evaluates each layer as an individual representation using a distance-weighted kNN criterion on validation data, and ranks layers by validation macro-F1 with balanced-accuracy and accuracy tie-breakers. It then constructs ranked layer prefixes and searches a controlled set of distance-fusion weights and kNN neighborhood sizes using validation data only. After selection, the layer subset, distance weights, and neighborhood size are frozen, and the selected configuration is evaluated once on the test split.

The method treats the internal hierarchy of a trained CNN as a set of candidate metric views and asks whether a compact, validation-selected fusion of layer-wise distances can improve classification over simpler alternatives. This framing is especially relevant in medical-image workflows, where reusing trained feature extractors can be attractive when labeled data, annotation expertise, computational cost, or deployment constraints limit full architecture redesign \cite{yosinski2014transferable,bar2024frofa,dorfner2025review}.

\method{} is evaluated through two complementary analyses. First, internal controls test whether the proposed design is needed. The native softmax head tests whether the fine-tuned CNN classifier is sufficient. Final-embedding classifiers test whether a classical classifier applied only to the deepest available representation is enough. Best-single-layer and uniform all-layer fusion variants then test two simpler explanations for any gain: that performance can be explained by one strong internal layer alone, or that it is sufficient to combine all available layers without validation-guided selection and weighting.

Second, external reference methods test whether established non-proposed mechanisms for reusing or combining embeddings already capture the same benefit. Multilayer and multiview implementations evaluate alternative ways of combining internal CNN views \cite{fradi2021multilayer,gomes2023fmlf,evci2022head2toe,aiolli2015easymkl,carroll1968gca,kettenring1971canonical,sharma2012gma,kan2016mvda}. Post-hoc metric and prototype baselines further test whether the observed behavior can be attributed to general embedding-space reuse rather than to the proposed validation-controlled multilayer distance-fusion rule \cite{goldberger2004nca,gautam2024kmex,boubekki2026posthoc}. This organization is important because a gain from internal embeddings may arise from several different mechanisms, including stronger local geometry, useful intermediate-depth information, reduced dependence on the final layer, multilayer complementarity, or model-selection effects. A fair evaluation must therefore distinguish between internal design evidence and comparison against external alternatives.

The contributions of this paper are:
\begin{itemize}
\item a validation-only algorithm for selecting internal CNN layers, distance-fusion weights, and kNN neighborhood size for brain tumor MRI classification;
\item a formulation of internal CNN reuse as a validation-controlled multilayer distance-fusion problem, separated from softmax classification, final-embedding classification, best-single-layer selection, blind all-layer fusion, and decision-level voting;
\item a design ablation showing that controlled multilayer aggregation improves on best-single-layer selection and uniform all-layer distance fusion on average;
\item a comparison against softmax heads, final-embedding classifiers, all-layer controls, multilayer fusion reference methods, multiview representation methods, and post-hoc metric or prototype baselines under the same train/validation/test protocol;
\item an aggregate analysis over 45 brain-MRI dataset-backbone contexts, with HAM10000 reported only as an external-domain control.
\end{itemize}
